% ============================================================================
%  presentation.tex  --  redesigned minimal academic Beamer theme
% ----------------------------------------------------------------------------
%  Keeps the bottom *progress bar* footline from the original "slate academic"
%  example (frame-aware, appendix-aware). Everything else is redesigned for a
%  cleaner, airier, more modern look.
% ============================================================================

\documentclass[aspectratio=169]{beamer}
\usepackage{pgf}

\title{Grant Selection Design and Start-Up Innovation}
\subtitle{Evidence from Design-Dependent Subsidy Effects}
\author{Felix W.\ Brandes}
\institute{Università degli Studi di Verona}
\date{June 2026}

% ---------------------------------------------------------------- palette ---
\definecolor{slate}{HTML}{3D5A80}      % primary structure (titles, frames) -- map blue
\definecolor{accent}{HTML}{E0683C}     % warm orange (Fig. 1 map) -> bar, rules, emphasis
\definecolor{teal}{HTML}{2E7D5B}       % emerald green (positive / good)
\definecolor{wash}{HTML}{EAF0F2}       % light background wash for blocks
\definecolor{ink}{HTML}{1C1C1C}        % body text
\definecolor{muted}{HTML}{7A8691}      % captions, subtitles, secondary text
\definecolor{danger}{HTML}{C44536}     % alerts / exclusion figures

% ------------------------------------------------------------- bemer colors -
\setbeamercolor{structure}{fg=slate}
\setbeamercolor{normal text}{fg=ink,bg=white}
\setbeamercolor{frametitle}{fg=slate,bg=white}
\setbeamercolor{framesubtitle}{fg=muted,bg=white}
\setbeamercolor{title}{fg=slate,bg=white}
\setbeamercolor{subtitle}{fg=muted,bg=white}
\setbeamercolor{author}{fg=ink,bg=white}
\setbeamercolor{date}{fg=muted,bg=white}
\setbeamercolor{institute}{fg=muted,bg=white}
\setbeamercolor{item}{fg=accent}
\setbeamercolor{subitem}{fg=accent!70}
\setbeamercolor{subsubitem}{fg=muted}
\setbeamercolor{block title}{fg=white,bg=slate}
\setbeamercolor{block body}{fg=ink,bg=wash}
\setbeamercolor{block title alerted}{fg=white,bg=danger}
\setbeamercolor{block body alerted}{fg=ink,bg=danger!8}
\setbeamercolor{block title example}{fg=white,bg=teal}
\setbeamercolor{block body example}{fg=ink,bg=teal!8}
\setbeamercolor{alerted text}{fg=danger}
\setbeamercolor{footline}{fg=muted,bg=white}
\setbeamercolor{caption name}{fg=slate}

% -------------------------------------------------------------- bemer fonts -
\setbeamerfont{frametitle}{size=\large,series=\bfseries}
\setbeamerfont{framesubtitle}{size=\small,series=\mdseries}
\setbeamerfont{title}{size=\LARGE,series=\bfseries}
\setbeamerfont{subtitle}{size=\normalsize,series=\mdseries}
\setbeamerfont{author}{size=\small,series=\mdseries}
\setbeamerfont{date}{size=\small,series=\mdseries}
\setbeamerfont{footline}{size=\scriptsize,series=\mdseries}
\setbeamerfont{block title}{size=\normalsize,series=\bfseries}
\setbeamerfont{caption}{size=\scriptsize}
\setbeamerfont{caption name}{size=\scriptsize,series=\bfseries}

\setbeamertemplate{navigation symbols}{}
\setbeamertemplate{caption}[numbered]
\setbeamertemplate{blocks}[rounded][shadow=false]
\pdfstringdefDisableCommands{\def\translate#1{#1}}

% --------------------------------------------------------- itemize markers --
\setbeamertemplate{itemize item}{%
  \usebeamercolor[fg]{item}\small$\blacktriangleright$%
}
\setbeamertemplate{itemize subitem}{%
  \usebeamercolor[fg]{subitem}\footnotesize$\triangleright$%
}
\setbeamertemplate{itemize subsubitem}{%
  \usebeamercolor[fg]{subsubitem}\scriptsize$\bullet$%
}
\setlength{\leftmargini}{1.4em}
\setbeamertemplate{itemize/enumerate body begin}{\vspace{0.3ex}}
% a little air between top-level items
\let\olditemize\itemize
\renewcommand{\itemize}{\olditemize\itemsep0.5ex\relax}

% ----------------------------------------------------------- frame title ----
\makeatletter
\setbeamertemplate{frametitle}{%
  \nointerlineskip%
  \vspace{0.45ex}%
  \begin{beamercolorbox}[wd=\paperwidth,sep=0pt]{frametitle}
    \hspace*{0.7cm}%
    \begin{minipage}{\dimexpr\paperwidth-1.4cm\relax}
      {\usebeamerfont{frametitle}\strut\insertframetitle\par}%
      \vspace{-0.15ex}%
      {\usebeamerfont{framesubtitle}\usebeamercolor[fg]{framesubtitle}%
        \strut\ifx\insertframesubtitle\@empty\else\insertframesubtitle\fi\par}%
      \vspace{-0.15ex}%
      % short, left-aligned accent rule (airier than a full-width line)
      {\color{accent}\rule{2.2cm}{1.4pt}}%
    \end{minipage}%
  \end{beamercolorbox}%
}
\addtobeamertemplate{frametitle}{\vspace*{1mm}}{\vspace*{0.2mm}}

% -------------------------------------------------------------- title page --
\defbeamertemplate*{title page}{slate modern}{%
  \vbox{}\vfill
  % cited-paper kicker
  \ifx\beamer@shortauthor\@empty\else\fi
  {\usebeamerfont{subtitle}\usebeamercolor[fg]{subtitle}%
    \scshape based on the paper\par}%
  \vspace{0.8ex}%
  \begin{beamercolorbox}[wd=\textwidth,leftskip=0pt]{title}
    \usebeamerfont{title}\inserttitle\par%
  \end{beamercolorbox}%
  \vspace{0.9ex}%
  {\color{accent}\rule{0.45\paperwidth}{1.6pt}}\par%
  \vspace{1.3ex}%
  \begin{beamercolorbox}[wd=\textwidth,leftskip=0pt]{subtitle}
    \usebeamerfont{subtitle}\insertsubtitle\par%
  \end{beamercolorbox}%
  \vfill
  \noindent
  \begin{minipage}[t]{0.62\textwidth}
    \raggedright
    {\usebeamerfont{author}\usebeamercolor[fg]{author}\insertauthor\par}%
    \vspace{0.6ex}%
    {\usebeamerfont{institute}\usebeamercolor[fg]{institute}\insertinstitute\par}%
  \end{minipage}\hfill
  \begin{minipage}[t]{0.34\textwidth}
    \raggedleft
    {\usebeamerfont{date}\usebeamercolor[fg]{date}\insertdate\par}%
  \end{minipage}%
  \vspace{1.6ex}%
}
\setbeamertemplate{title page}[slate modern]

% =============================================================== PROGRESS BAR
%  (preserved from the user's original example, lightly cleaned)
% ----------------------------------------------------------------------------
\newlength{\beamerstyle@progresswidth}
\newif\ifbeamerstyle@appendix
\beamerstyle@appendixfalse
\let\beamerstyle@oldappendix\appendix
\renewcommand{\appendix}{%
  \beamerstyle@appendixtrue%
  \beamerstyle@oldappendix%
  \setcounter{framenumber}{0}%
}

\newcommand{\beamerstyle@setmainframes}{%
  \edef\beamerstyle@mainframes{\inserttotalframenumber}%
  \ifcsname insertmainframenumber\endcsname%
    \ifx\insertmainframenumber\@empty\else%
      \edef\beamerstyle@mainframes{\insertmainframenumber}%
    \fi%
  \fi%
  \edef\beamerstyle@displayframe{\insertframenumber}%
  \ifbeamerstyle@appendix%
    \edef\beamerstyle@mainframes{\inserttotalframenumber}%
    \ifcsname appendixtotalframenumber\endcsname%
      \edef\beamerstyle@mainframes{\appendixtotalframenumber}%
    \fi%
    \@ifundefined{beamerstyle@appendixtotal}{}{%
      \edef\beamerstyle@mainframes{\beamerstyle@appendixtotal}%
    }%
  \fi%
}

\setbeamertemplate{footline}{%
  \leavevmode%
  \begingroup%
  \beamerstyle@setmainframes%
  \pgfmathsetmacro{\beamerstyle@progressratio}{%
    min(\beamerstyle@displayframe/max(\beamerstyle@mainframes,1),1)%
  }%
  \setlength{\beamerstyle@progresswidth}{\beamerstyle@progressratio\paperwidth}%
  \makebox[0pt][l]{\color{muted!28}\rule{\paperwidth}{0.55pt}}%
  \makebox[0pt][l]{\color{accent}\rule{\beamerstyle@progresswidth}{0.9pt}}%
  \par\vspace{1.0ex}%
  \hbox to \paperwidth{%
    \hfill%
    \usebeamerfont{footline}\usebeamercolor[fg]{footline}%
    \beamerstyle@displayframe\,/\,\beamerstyle@mainframes\hspace{0.7cm}%
  }%
  \vspace{0.6ex}%
  \endgroup%
}
\makeatother

% No persistent header breadcrumb.
\setbeamertemplate{headline}{}

% ===================================================== speaker-notes page ====
% Note page: section breadcrumb (accent mark, top-left) + small slide thumbnail
% (top-right) + the note text. NO date, no grey spine.
\makeatletter
\defbeamertemplate{note page}{clean}{%
  \nointerlineskip
  \begin{minipage}[t]{0.60\linewidth}\vspace{0pt}%
    {\footnotesize\color{muted}\insertshorttitle}\par\vspace{1pt}%
    {\color{accent}\rule[-0.35ex]{1.1mm}{2.7ex}}\hspace{2mm}%
    {\normalsize\color{slate}\bfseries
       \ifbeamerstyle@appendix Backup / Q\&A\else\insertsection\fi}%
  \end{minipage}\hfill
  \begin{minipage}[t]{0.37\linewidth}\vspace{0pt}\raggedleft
    \insertslideintonotes{0.30}%
  \end{minipage}\par
  \vspace{1.5mm}{\color{accent!45}\rule{\linewidth}{0.7pt}}\par\vspace{3mm}%
  \setbeamertemplate{itemize item}{\small\color{accent}\textbullet}%
  \setbeamertemplate{itemize subitem}{\scriptsize\color{accent!70}--}%
  \setbeamerfont{itemize/enumerate body}{size=\small}%
  \setlength{\leftmargini}{1.5em}%
  {\small\insertnote}%
}
\makeatother
\setbeamertemplate{note page}[clean]

% ===================================================== section divider slides
\setbeamertemplate{section in toc}{%
  \leavevmode\leftskip=0pt%
  {\color{accent}$\blacktriangleright$}\,\inserttocsection\par%
}
\setbeamertemplate{subsection in toc}{}  % keep dividers clean

\makeatletter
\AtBeginSection[]{%
  \ifbeamerstyle@appendix\else
    \begin{frame}[plain,noframenumbering]
      \vfill
      {\color{accent}\rule{0.16\paperwidth}{2.4pt}}\par
      \vspace{1.0ex}
      {\usebeamerfont{title}\usebeamercolor[fg]{title}\insertsection\par}
      \vspace{1.4ex}
      {\scriptsize
        \begin{minipage}{0.82\paperwidth}
          \setlength{\parskip}{0.2ex}
          \tableofcontents[currentsection,hideallsubsections,%
                           sectionstyle=show/shaded]
        \end{minipage}}
      \vfill
    \end{frame}%
  \fi
}
\makeatother


\begin{document}

\begin{frame}[plain,noframenumbering]
  \titlepage
\end{frame}

% ---------------------------------------------------------------- Slide 1 ---
\begin{frame}{Curriculum Vitae}
  \begin{itemize}
    \item \textbf{School} \\
      Sidcot Independent School {\color{muted}\small England \hfill 2012--2014} \\
      {\color{muted}\small International Baccalaureate}
    \vspace{0.8ex}
    \item \textbf{University} \\
      Technical University of Munich {\color{muted}\small B.Sc. Technology and Management \hfill 2014--2020} \\  
      Technical University of Munich {\color{muted}\small M.Sc. Technology and Management \hfill 2020--2024} \\ 
    \vspace{0.8ex}
    \item \textbf{Professional Experience} \\
      CoInvest Corporate Finance {\color{muted}\small Analyst \hfill 2021--2024} \\  
      agentix GmbH {\color{muted}\small Co-Founder \& COO \hfill 2024--2025} \\
      Independent Management Consultant {\color{muted}\small  \hfill 2024--present}
  \end{itemize}
\end{frame}

% ---------------------------------------------------------------- Slide 2 ---
\begin{frame}{Master Thesis}
  \begin{block}{Research Question}
    How can you improve investments into Start-Ups and SMEs through
    legislative change of Investment Modalities?
  \end{block}
\end{frame}

% ---------------------------------------------------------------- Slide 3 ---
\begin{frame}{Agentix}
  \begin{itemize}
    \item AI at the forefront of Innovation
    \item \textit{``Automating Cognitive Workflows''}
    \item Recipient of EXIST Grant
  \end{itemize}
\end{frame}

% ---------------------------------------------------------------- Slide 4 ---
\begin{frame}{Research Proposal: Overview}
    \begin{block}{Core Question} Do the selection signals used in public SME and start-up funding programmes generate technological innovation, observable spillover, and economic value? 
    \end{block}
  \vspace{0.6ex}
  \begin{itemize}
    \item \textbf{Motivation:} Private returns to innovation fall short of social returns
      (Arrow, 1962) --- the appropriability problem is sharpest for early-stage firms
    \item \textbf{Gap:} Literature evaluates \emph{whether} grants work; not \emph{which
      selection signals} generate which outcomes
    \item \textbf{Contribution:} Shifts programme evaluation from beneficiary status to the
      economic content of the selection rules
  \end{itemize}
\end{frame}

% ---------------------------------------------------------------- Slide 5 ---
\begin{frame}{Theoretical Framework}
  \textbf{Two-step design}
  \vspace{0.5ex}
  \begin{itemize}
    \item[\color{accent}1.] \textbf{Identify the selection rule} --- code programme documents,
      scoring rubrics and eligibility criteria; recover implied signal weights; compare
      stated vs.\ applied weights
    \item[\color{accent}2.] \textbf{Estimate causal effects} --- regression-discontinuity around
      score / budget cut-offs; applicants just above and below the threshold are otherwise
      comparable
  \end{itemize}
  \vspace{0.8ex}
  \textbf{Three outcome layers}
  \vspace{0.3ex}
  \begin{itemize}
    \item \textbf{Technological Innovation} --- patents, product/process launches, technology-class novelty
    \item \textbf{Observable Spillover} --- licensing, research collaborations, technology diffusion
    \item \textbf{Economic Value} --- employment, enterprise survival, cluster/sector growth
  \end{itemize}
\end{frame}

\begin{frame}{Empirical Framework} 
\begin{itemize} 
  \item \textbf{Unit of analysis:} enterprise--application pair. 
  \item \textbf{Object of interest:} observable application signals used in the funding decision. 
  \item \textbf{Treatment margin:} applicants just above and just below score or budget cut-offs. 
  \item \textbf{Preferred design:} regression discontinuity around selection thresholds. 
  \item \textbf{Outcome hierarchy:} 
  \begin{itemize} 
    \item technological innovation
    \item observable spillover
    \item local economic value
  \end{itemize} 
\end{itemize} 
\end{frame}


% ---------------------------------------------------------------- Slide 6 ---
\begin{frame}{Research Proposal: Data and Methodology}
  \begin{columns}[T]
    \begin{column}{0.48\textwidth}
      \textbf{Data Sources}
      \begin{itemize}
        \item Programme documents, scoring rubrics, beneficiary lists (EU transparency mandate)
        \item CORDA / ORBIS --- applicant-level award data
        \item EPO / Patstat --- patent records
        \item Business registers \& regional employment data
        \item EIC Accelerator score-level data
      \end{itemize}
    \end{column}
    \begin{column}{0.48\textwidth}
      \textbf{Identification Strategy}
      \begin{itemize}
        \item \textbf{Preferred:} Regression discontinuity at score / budget cut-offs
          (Imbens \& Lemieux, 2008)
        \item \textbf{Alternative:} DiD / event study across calls, regions, countries
          (Callaway \& Sant'Anna, 2021)
        \item \textbf{Threats:} Self-selection into application, score manipulation near cut-offs,
          survivorship bias
      \end{itemize} 
    \end{column}
  \end{columns}
\end{frame}



\begin{frame}{Expected Contribution} 
\begin{itemize} 
  \item \textbf{Empirical:} estimate whether selection signals generate the outcomes used to justify public support. 
  \item \textbf{Conceptual:} separate the effect of receiving a grant from the design of the rules that allocate grants. 
  \item \textbf{Policy-relevant:} identify where programme objectives and actual selection practice diverge.
  \item \textbf{PhD logic:} the project can begin as one empirical paper on grant selection design and develop into a broader research agenda on innovation policy. 
\end{itemize} 
\end{frame}

\begin{frame}{Backup: Data Sources}
  \begin{itemize}
    \item Programme documents, eligibility criteria, scoring rubrics and beneficiary lists.
    \item Applicant-level award data from CORDA / ORBIS where accessible.
    \item Patent data from EPO / PATSTAT.
    \item Firm accounts, business registers and regional employment data.
    \item Score-level data from programmes such as the EIC Accelerator where available.
  \end{itemize}
\end{frame}

\begin{frame}{Backup: Identification Risks}
  \begin{itemize}
    \item Self-selection into application.
    \item Score manipulation near cut-offs.
    \item Missing or incomplete score-level data.
    \item Survivorship bias in firm-level outcomes.
    \item Outcome correlation across innovation, spillover and economic-value layers.
  \end{itemize}
\end{frame}

\end{document}
